\begin{pastexam}[label={px:ntu-114-bime-ds-02}]{臺大生機所 114 學年度 第 2 題}
（虛擬碼；10 分）

\textbf{(10\%)} Write a recursive function in pseudocode that takes a reference (or pointer) to the
first node of a linked list as argument and returns the value of the maximum key in the list.
Assume that the keys in the linked list are all positive integers. If the linked list is empty, the
function should return 0.
\end{pastexam}

\begin{pastsoln}
節點有欄位 $\mathit{key}$ 和 $\mathit{next}$（資結書第~5.1~節），$x$ 指向第一個節點，空串列是
$x = \NIL$。

遞迴的想法：串列的最大值是「第一個節點的鍵值」和「其餘部分的最大值」兩者中較大的。 其餘部分
是以 $x.\mathit{next}$ 開頭的串列，長度少一，交給遞迴。 空串列回傳 $0$；鍵值都是正整數，所以
$0$ 比任何鍵值都小，拿來和第一個鍵值比不會出錯，最後一個節點也不用特別處理。

\begin{pseudo}{遞迴求串列的最大鍵值}
\Function{List-Max}{$x$}
    \If{$x = \NIL$}
        \State \Return $0$ \Comment{空串列}
    \EndIf
    \State $m \gets$ \Call{List-Max}{$x.\mathit{next}$} \Comment{其餘部分的最大值}
    \If{$x.\mathit{key} > m$}
        \State \Return $x.\mathit{key}$
    \Else
        \State \Return $m$
    \EndIf
\EndFunction
\end{pseudo}

例：串列 $3 \to 9 \to 4$。 \proc{List-Max}$(\NIL) = 0$；節點 $4$：$\max(4, 0) = 4$；節點 $9$：
$\max(9, 4) = 9$；節點 $3$：$\max(3, 9) = 9$，回傳 $9$。

\textbf{正確性。} 對串列長度 $L$ 歸納。 $L = 0$：回傳 $0$，符合題目要求。 $L \ge 1$：由歸納假設
$m$ 是後面 $L - 1$ 個鍵值的最大值（$L = 1$ 時 $m = 0$，小於任何正整數鍵值），函式回傳
$\max(x.\mathit{key}, m)$，正是全部 $L$ 個鍵值的最大值。

\textbf{複雜度。} 每個節點呼叫一次，加上最後的 \NIL 一次，共 $L + 1$ 次呼叫，每次 $O(1)$，時間
$\Theta(L)$。 遞迴深度 $L + 1$，呼叫堆疊用 $\Theta(L)$ 空間；同樣的事用迴圈寫只要 $O(1)$ 空間，
但題目要求遞迴。
\end{pastsoln}
